\hwhat{The units are per-agent trap spells: inactive spells, timer-pause chains, error loops and repeated-command loops. We test them in 11 non-holdout goal periods. Kramers escape predicts a constant escape hazard. Instead, traps \emph{age}: within an agent, the hazard falls as about $t^{-0.8}$ (\#51 $\beta=-0.77$, \#38 $-1.71$). Each re-pause lowers the odds of action at the next timer expiry (\#51 $-0.38$). One directed message during a pause raises those odds $\times$1.5--2.9 (5/7 periods); more messages add little. Round 1b counts real failures instead of stderr, takes the leading-@ nudge target and adds Jev blocked and loop spells. Error-loop aging was mostly a stderr artifact: loop rows fall 51\%, and aging survives only in \#51 ($-0.57$). Jev blocked spells age in \#51 and \#27. The kick results do not change. With the nudger off (NE43), aging does not change. Mean-field $\beta J_0=0.09$--$0.41$ gives no bistability. The holdout is not run.}

\hmodel{Models 01/02 (mean field) with Kramers escape. The activity $x$ of an agent moves in a double well $U(x)$ (idle and active basins) at temperature $T$. A trap is a stay in the idle basin. Escape has rate $k$ over barrier $\Delta U$, with attempt rate $\nu$. A kick (a message that reaches the agent) adds a force $f$. Swarm: Curie--Weiss coupling $J_0$, field $h$, active fraction $m$. Fit: hazard $h_{it}$, agent effects $\alpha_i$, aging slope $\beta$, dose effects $\gamma_n$.}

\hmath{\vspace{-5pt}\begin{gather*}
k=\nu\,e^{-\Delta U/T},\qquad G(x)=-\ln P(x)=U(x)/T+\text{const}\\
k(n)=k_0\,e^{nf/T}\ \text{(Kramers)}\quad\text{vs}\quad k(n)=k_0(1+c\,n)\ \text{(triggers)}\\
\ln\!\big[-\ln(1-h_{it})\big]=\alpha_i+\beta\ln t+\textstyle\sum_{n}\gamma_n\,\mathbb 1[\text{dose}_{it}=n]\\
f(m)=-\tfrac{1}{2}J_0m^2-hm-Ts(m);\ \text{bistable iff }\beta J_0(1-m^2)>1
\end{gather*}\vspace{-7pt}}

\hobs{../figures/summary_obs.pdf}{(a) Escape hazard from inactive spells vs time in the spell, in the two periods with the most power. A spell ignores single stray actions. After the first minutes (grey), Kramers predicts the flat dashed line. The measured hazard falls, so traps age. (b) \#51 hazard and odds ratios vs the number of kicks. Dotted lines extrapolate Kramers (exponential in dose) from one kick. The measured curves flatten: the first message carries the effect.}

\htelos{Kick a stuck agent early: its escape hazard falls as about $t^{-0.8}$ with time in the trap. Send one directed message during a pause. It raises the odds of action at the next timer expiry $\times$1.5--2.9; repeats and room messages add little. Under the 5-min pause default (after NE44), a kick lifts escape from deep pause chains only from 0.17 to 0.25. Messages do not break real failure loops (9/9 periods), so use tooling or a reset there. The aging is the agent's own: it does not change with the nudger off. The physics adds little, because Kramers fails and aging with saturating triggers is ordinary survival analysis.}

\hbottom{Traps age by themselves, not by Kramers escape; one directed kick at a pause gate helps, and real failure loops age only in \#51.}

\hresults{\footnotesize\begin{tabular}{@{}p{0.34\columnwidth}p{0.44\columnwidth}l@{}}
\hline
prediction (round 1b) & observed & verdict\\\hline
real-failure loops age (majority) & 3/10 point, 0/10 CI; \#51 $-0.57$ & failed\\
messages do not break them & 9/9 periods & holds\\
Jev blocked spells age (\#51, \#38, $\ge$half) & \#51 $-0.39$, \#27 $-1.24$; \#38 flat; timer-like in 3 & failed\\
Jev loop spells age (\#51) & $-0.46$ [$-0.55,-0.18$] & partial\\
gate kick OR $\ge1.5$ (leading @) & 5/7 (1.5--2.9) & holds\\
room kicks above null & 6/8 & holds\\
NE44 kick$\times\ln k$: 0 then $<0$ & $+0.57$ before, $+0.06$ after & failed\\
NE44 declared pauses shrink & 90 s $\to$ 180 s & failed\\
NE43 aging without kicks & TS1r $-0.59\to-0.53$; gates $-0.31\to-0.26$ & holds\\
NE43 gate escape drops & $+0.006$ [$-0.09,0.10$] & failed\\
\hline
\end{tabular}}

\hobsb{../figures/r1b_summary.pdf}{(a) Loop aging slope per period: error loops on stderr (round 1), real failures (round 1b) and Jev blocked spells. The gray band is memoryless. Only \#51 keeps a clear negative slope on real failures. (b) Gate escape vs pause-chain depth. Before NE44, a directed kick rescues agents at any depth. After NE44, deep chains barely respond.}

\hcaveats{The NE44 contrast compares periods, because the change sits inside the locked holdout window. Day length (4 h vs 8 h), roster and goals confound it. Jev \texttt{p\_blocked} means ``waiting on others'' in regime I and ``failing or looping'' in regime III. Window spells are short in 4-h periods. Kick labels are temporal coincidences. Several 3--5-day periods have unstable day-bootstrap CIs; we report Wald CIs alongside. We did not re-run the barrier-inversion and swarm blocks, because their inputs did not change.}

\hnext{Test self-reinforcement with a Pólya urn on the context ledger: own-output share vs the aging slope. Test content-level spirals on the holdout (\#45).}
